% -----------------------------------------------------------------------------------------------------
\section{$\MSbar$ 部分子分布的正定性}\label{sec:front-positivity}

裸的光前数密度诠释提示 $f\ge0$，但重整化（对算符的紫外发散作 $\MSbar$ 减除）破坏了任何自动的正定性。2020--2024 年的争论澄清了下面两点：$\MSbar$ 分布确实可以为负，而在足够大的标度上为正只有部分论证（核心陈述\cref{conj:pdf-positivity}，[C]）。

\begin{proposition}[$\MSbar$ 部分子分布可以为负 \tagV]\label{prop:front-msbar-negative}
(a)（单圈 Yukawa 模型，\tagR，$R_{\mathrm{PT}}$）（$R_{\mathrm{PT}}$：在重整化微扰论中对红外正规化的部分子态严格。）在 Yukawa 理论 $\mathcal L_{\rm int}=-\lambda\bar\Psi_N\psi_q\phi+{\rm h.c.}$ 中（全部粒子有质量，弱耦合），“核子”$\Psi_N$（质量 $m_p$）中夸克（$m_q$）的单圈 $\MSbar$ 分布为
\begin{equation}\label{eq:front-yukawa}
\begin{gathered}
f^{\rm renorm}_{q/p}(\xi;\mu)=a_\lambda(\mu)\,(1-\xi)\Bigl[\ln\frac{\mu^2}{\Delta(\xi)}-1+\frac{(m_q+\xi m_p)^2}{\Delta(\xi)}\Bigr]+\order{a_\lambda^2},\\
\Delta(\xi)=\xi m_s^2+(1-\xi)m_q^2-\xi(1-\xi)m_p^2,
\end{gathered}
\end{equation}
$a_\lambda=\lambda^2/(16\pi^2)$，$m_s$ 为标量质量；对稳定的靶（$m_p<m_q+m_s$）$\Delta>0$。对每个 $\xi$，它在 $\mu$ 足够小时为负、足够大时为正。取 $m_q=0.3$、$m_p=1.0$、$m_s=1.5\,\mathrm{GeV}$：$\mu=0.8\,\mathrm{GeV}$ 时对 $\xi>0.264$ 为负，$\mu=1.2\,\mathrm{GeV}$ 时对 $\xi>0.567$ 为负，$\mu=2.5\,\mathrm{GeV}$ 时对一切 $\xi$ 为正；$\xi\to1$ 处的变号点为 $\mu_*=m_s\exp\{\frac12[1-(m_q+m_p)^2/m_s^2]\}=1.70\,\mathrm{GeV}$，$\xi\to0$ 时方括号趋于 $\ln(\mu^2/m_q^2)$。
(b) 因此“$\MSbar$ 部分子分布在所有标度上非负”作为一般定理不成立。（本命题是 Yukawa 模型中的微扰陈述。）
% VF-FRO-10 (exact thresholds 0.264, 0.567; WRITER_SPEC_v2 6.11 item 7(c))
\uses{def:pdf-quark, def:prelim-msbar}
\end{proposition}

\begin{proof}[(a) 的证明]
按\cref{def:pdf-quark}（没有规范场，故没有 Wilson 线），最低阶割图给出
\begin{equation*}
f(\xi)=\frac12\int\frac{\dd k^-\dd^2k_T}{(2\pi)^4}\Tr[\gamma^+\Phi(k)],\qquad
\Phi=2\pi\delta\bigl((P-k)^2-m_s^2\bigr)\,\lambda^2\,\frac{(\sla k+m_q)\frac12(\sla P+m_p)(\sla k+m_q)}{(k^2-m_q^2)^2},
\end{equation*}
$k^+=\xi P^+$（对靶自旋平均）。在割线上 $\Tr[\gamma^+(\sla k+m_q)(\sla P+m_p)(\sla k+m_q)]=4P^+[k_T^2+(m_q+\xi m_p)^2]$，$k^2-m_q^2=-[k_T^2+\Delta(\xi)]/(1-\xi)$，$\delta((P-k)^2-m_s^2)=\delta(k^--k^-_0)/[2(1-\xi)P^+]$，于是
\begin{equation*}
f(\xi)=\frac{\lambda^2}{16\pi^3}(1-\xi)\int\dd^2k_T\,\frac{k_T^2+(m_q+\xi m_p)^2}{[k_T^2+\Delta(\xi)]^2}.
\end{equation*}
在 $D=4-2\varepsilon$ 维中用 $\MSbar$ 耦合（\cref{def:prelim-msbar}），横向测度给出 $a_\lambda(1-\xi)\frac{(\mu^2\ee^{\gamma_E})^\varepsilon}{\Gamma(1-\varepsilon)}\int_0^\infty\dd k\,k^{-\varepsilon}\frac{k+A}{(k+\Delta)^2}$，$A=(m_q+\xi m_p)^2$。由 $\int_0^\infty\dd k\,\frac{k^{1-\varepsilon}}{(k+\Delta)^2}=\Delta^{-\varepsilon}\Gamma(2-\varepsilon)\Gamma(\varepsilon)$ 与 $\int_0^\infty\dd k\,\frac{k^{-\varepsilon}}{(k+\Delta)^2}=\Delta^{-1-\varepsilon}\Gamma(1-\varepsilon)\Gamma(1+\varepsilon)$，裸分布为 $a_\lambda(1-\xi)[\frac1\varepsilon+\ln\frac{\mu^2}{\Delta}-1+\frac A\Delta]+\order{\varepsilon}$；$\MSbar$ 抵消项 $-a_\lambda(1-\xi)/\varepsilon$ 即得\cref{eq:front-yukawa}（极点的系数 $1-\xi$ 是该模型的单圈劈裂函数）。$\Delta=\xi(1-\xi)[\frac{m_q^2}\xi+\frac{m_s^2}{1-\xi}-m_p^2]$，而 $\frac{m_q^2}\xi+\frac{m_s^2}{1-\xi}-(m_q+m_s)^2=\frac{[m_q(1-\xi)-m_s\xi]^2}{\xi(1-\xi)}\ge0$，故 $m_p<m_q+m_s$ 时 $\Delta>0$。方括号对 $\ln\mu^2$ 的导数为 $1$，在 $\mu\to0$ 时趋于 $-\infty$、$\mu\to\infty$ 时趋于 $+\infty$，所以对每个 $\xi$ 恰在 $\mu_0(\xi)=\sqrt{\Delta}\,\ee^{(1-A/\Delta)/2}$ 处变号；$\xi\to1$ 时 $\Delta\to m_s^2$、$A\to(m_q+m_p)^2$，给出 $\mu_*$；$\xi\to0$ 时 $\Delta\to m_q^2$、$A/\Delta\to1$。所列数值由求方括号的零点得到；数值上 $\mu_0(\xi)$ 从 $m_q$ 单调增到 $\mu_*$，故对一切 $\xi$ 为正当且仅当 $\mu\ge\mu_*$。最后，在 $k_T<\mu$ 处截断的分布 $f_{\rm cut}$ 处处为正，并且精确地有 $f^{\rm renorm}=f_{\rm cut}-a_\lambda(1-\xi)[\ln(1+\Delta/\mu^2)-(A-\Delta)/(\mu^2+\Delta)]$：两者只差 $\order{m^2/\mu^2}$，而在 $\mu\lesssim m$ 时 $\MSbar$ 的渐近 $1/k_T^2$ 减除“减多了”（这正是\cite[第~VIII.B~节]{CollinsRogersSato2022}的机制）。
% checks: writing/work/W27/s_yukawa.py (trace, DR, MS-bar, zeros 0.26403/0.56737, mu* = 1.6988, monotonic mu0) and
% s_yukawa_cutoff.py (exact cut-off relation); frontier check OTHERS s2
\end{proof}

\paragraph{地位与缺口。}[V]，证据：显式的微扰反例（模型），即\cite{CollinsRogersSato2022}式~(55)--(59) 的模型。分歧（三方并列，不作裁决）：\cite{CandidoForteHekhorn2020}（据摘要）论证 NLO $\MSbar$ PDF “在微扰区”为正（先构造正的因子化方案，再在 NLO 的方案变换小于领头阶时变换到 $\MSbar$）；\cite{CollinsRogersSato2022}指出该论证在“裸 PDF 为正”一步失效：在他们称为 track B 的做法（无质量部分子量的共线发散用 $n>4$ 正规化，并与裸强子 PDF 卷积）中，紫外发散贡献在 $n<4$ 时为正，但“can and often do become negative after analytic continuation to $n>4$”，而在 track A（强子的重整化算符 PDF，部分子计算只用于硬系数）中，定义与因子化定理都不排除低、中等标度上为负的 $\MSbar$ PDF；\cite[第~3--4~节]{Candido:2023ujx}则认为 PDF “always turn positive at high enough scale”且“turn negative at a low enough scale”，后者是非微扰标度。一个物理例子：$\MSbar$ 重夸克 PDF 在匹配阈值稍下方为负\cite[第~I~节]{CollinsRogersSato2022}。误解 M2：“PDF 是概率密度，所以 $\MSbar$ PDF 为正”。

%%FIGBLOCK fig:front-yukawa-pdf%%

\begin{example}[Yukawa 模型分布的端点行为与符号变化]\label{ex:front-yukawa}
$\xi\to0$：$\Delta\to m_q^2$，$(m_q+\xi m_p)^2/\Delta\to1$，方括号 $\to\ln(\mu^2/m_q^2)$，当且仅当 $\mu<m_q$ 时为负。$\xi\to1$：$\Delta\to m_s^2$，方括号 $\to\ln(\mu^2/m_s^2)-1+(m_q+m_p)^2/m_s^2$，在 $\mu_*^2=m_s^2\ee^{1-(m_q+m_p)^2/m_s^2}$ 处为零；对上述质量，$\mu_*=1.5\,\ee^{\frac12(1-1.69/2.25)}\,\mathrm{GeV}=1.70\,\mathrm{GeV}$。\cref{fig:front-yukawa-pdf}画出三个标度上的曲线：$\mu=0.8\,\mathrm{GeV}$ 的曲线在 $\xi\approx0.58$ 处取最小值约 $-0.35\,a_\lambda$。
\uses{prop:front-msbar-negative}
\end{example}

\begin{proposition}[足够大标度上的正定性 \tagP]\label{prop:front-positivity-large-scale}
设 (PS1) 一组正的观测量（如 DIS 结构函数，\cref{thm:pdf-dis-xspace}）在领头幂次共线因子化，并定义一个“物理”方案；(PS2) 高扭度可忽略；(PS3) NLO 方案变换是微扰的（\cref{rem:pdf-schemes}）。则 $\MSbar$ PDF 在某个标度以上为正；由于对角道 $\order{\alpha_s}$ 项的 Sudakov 符号（$\MSbar$ 对共线奇异性减多了，而对角道中它们具有 Sudakov 符号），该标度在大 $x$ 处与 $x$ 无关；\cite{Candido:2023ujx}保守地估计为 $Q^2\gtrsim5\,\mathrm{GeV}^2$（来自 $x\sim0.3$、$\alpha_s\lesssim0.2$ 时 $\frac{\alpha_s}{2\pi}C_{gg}(x)<1$）。\cite{CollinsRogersSato2022}定性上同意：“once the $\overline{\rm MS}$ scale is large enough compared with intrinsic transverse momentum scales, the pdfs indeed become positive”。非微扰地位：\cref{asm:intro-wightman}、\cref{asm:intro-composite}、\cref{asm:intro-bridge}（MA0、MA1+MA2）。
\uses{thm:pdf-dis-xspace, rem:pdf-schemes}
\end{proposition}

\paragraph{地位与缺口。}[P]：在 (PS1)--(PS3) 之下的 NLO 论证，高阶与幂次修正不受控制，阈值标度是非微扰的。

\begin{remark}[对本书与拟合的后果]\label{rem:front-positivity-tracks}
(1) $\MSbar$ PDF 的正定性与 Soffer 型不等式在本书中从不作为定理陈述（\cref{prop:coll-soffer}只断言领头阶演化保持后者）；(2) 在低输入标度上把正定性作为拟合约束是一种理论偏置\cite{CollinsRogersSato2022}；(3) 固定阶因子化截面可以为负而不矛盾\cite[附录~A]{CollinsRogersSato2022}；(4) 正的方案（$k_T$ 截断、物理方案）存在——在模型中单圈显式成立（\cref{prop:front-msbar-negative}的证明），其他情形只有一般论证。
\end{remark}
